% SGA 3, Expose IV, all of section 4.7.
% French authority: Exp4-13oct24.pdf, local pages 53--55
% (combined-reader pages 241--243; printed pages 229--231).
% The external English Markdown was used only as a comparison.

\setcounter{footnote}{45}

\SGARefTarget{sga3:iv:subsection:4:7}\subsection*{4.7. 有效性判别法的应用: 同构定理}
\addcontentsline{toc}{subsection}{\SGARefLink{sga3:iv:subsection:4:7}{4.7}. 有效性判别法的应用: 同构定理}

本小节给出一个有效性判别法的应用实例.
初始数据为: \(\mathscr C\) 上的一个拓扑 \(\mathcal T\),
仍假设它比典范拓扑更粗;
\(\mathscr C\) 的一个态射族 \((M)\),
满足 \SGARefLink{sga3:iv:para:4:6:3}{4.6.3} 的公理
\SGARefLink{sga3:iv:para:4:6:3:item:a:01}{(a)}--\SGARefLink{sga3:iv:para:4:6:3:item:e:mathcal:t:01}{\((e_{\mathcal T})\)};
以及 \(\mathscr C\) 的一个态射族 \((N)\), 它可以满足以下公理:
\begin{enumerate}
  \item[(a)]\SGARefTarget{sga3:iv:subsection:4:7:item:a:01}
  \((N)\) 在基变换下稳定;
  \item[\((f_{\mathcal T})\)]\SGARefTarget{sga3:iv:subsection:4:7:item:f:mathcal:t:01}
  \((N)\) 中的态射关于给定拓扑可以下降.
  更确切地说, 对每个 \(S\in\operatorname{Ob}(\mathscr C)\),
  每个覆盖族 \(\{S_i\to S\}\), 以及每个属于 \((N)\) 的态射族
  \(\{T_i\to S_i\}\),
  \(T_i\) 上关于 \(\{S_i\to S\}\) 的每个下降数据都有效;
  若 \(T\) 表示下降所得的对象, 则 \(T\to S\) 属于 \((N)\).
\end{enumerate}

由于 \((M)\) 的每个元素都是覆盖态射
(条件 \SGARefLink{sga3:iv:para:4:6:3}{4.6.3}
\SGARefLink{sga3:iv:para:4:6:3:item:d:mathcal:t:01}{\((d_{\mathcal T})\)}),
公理 \SGARefLink{sga3:iv:subsection:4:7:item:f:mathcal:t:01}{\((f_{\mathcal T})\)}
蕴含\SGARefLink{sga3:iv:subsection:4:7:item:f:m:01}{以下公理}:
\footnote{编者注: 公理
\SGARefLink{sga3:iv:subsection:4:7:item:f:m:01}{\((f_M)\)} 原置于命题
\SGARefLink{sga3:iv:proposition:4:7:2}{4.7.2} 之前, 现移至此处.}
\begin{enumerate}
  \item[\((f_M)\)]\SGARefTarget{sga3:iv:subsection:4:7:item:f:m:01}
  若 \(Y\to X\) 属于 \((N)\), 且 \(X\to X_1\) 属于 \((M)\),
  则 \(Y\) 上关于 \(X\to X_1\) 的每个下降数据都有效;
  若 \(Y_1\) 表示下降所得的对象, 则 \(Y_1\to X_1\) 属于 \((N)\).
\end{enumerate}

以下是一个将在后文处理的例子:
\(\mathscr C\) 是概形范畴,
\(\mathcal T\) 是忠实平坦拟紧拓扑,
\((M)\) 是忠实平坦拟紧态射族,
而 \((N)\) 是闭浸入族, 或拟紧浸入族.
\footnote{编者注: 参见 \SGARefLink{sga3:iv:subsection:6:4}{\S6.4};
另见 VI\(_A\), \SGARefLink{par:sga3-VIA-5.3.1}{5.3.1}.}

结合 \SGARefLink{sga3:iv:remark:4:6:11}{4.6.11},
可将 \SGARefLink{sga3:iv:theorem:4:6:6}{4.6.6} 的主要结果重述如下.

\medskip
\noindent\SGARefTarget{sga3:iv:proposition:4:7:1}\textbf{命题 4.7.1.}
\textit{设 \(X\) 是 \(\mathscr C\) 的一个可平方对象,
\(R\) 是 \(X\) 中一个 \((M)\) 型的等价关系.
假设 \(R\to X\times X\) 属于 \((N)\),
其中 \((N)\) 满足 \SGARefLink{sga3:iv:subsection:4:7:item:a:01}{(a)} 和
\SGARefLink{sga3:iv:subsection:4:7:item:f:mathcal:t:01}{\((f_{\mathcal T})\)}.
则层商 \(X/R\) 由下式给出:}
\[
\begin{aligned}
 (X/R)(S)=\{Z\subseteq X_S\mid{}&
   Z\text{ 是一个 }S\text{-子对象, 且在 }R\times S\text{ 下稳定},\\
 & Z\to X_S\text{ 属于 }(N),\quad
   Z\to S\text{ 是覆盖态射或属于 }(M),\\
 & R_Z\simeq Z\times_S Z\,\}.
\end{aligned}
\]

还有\SGARefLink{sga3:iv:proposition:4:7:2}{以下结果}.

\medskip
\noindent\SGARefTarget{sga3:iv:proposition:4:7:2}\textbf{命题 4.7.2.}
\textit{设 \(X\in\operatorname{Ob}(\mathscr C)\),
\(R\) 是 \(X\) 中的一个 \((M)\)-有效等价关系,
\((N)\) 是一个满足 \SGARefLink{sga3:iv:subsection:4:7:item:a:01}{(a)} 和
\SGARefLink{sga3:iv:subsection:4:7:item:f:m:01}{\((f_M)\)} 的态射族.
对 \(X\) 的每个在 \(R\) 下稳定且使得 \(Y\to X\) 属于 \((N)\)
的子对象 \(Y\), \(R\) 在 \(Y\) 中诱导的等价关系 \(R_Y\)
是 \((M)\)-有效的. 此外, 商}
\[
  Y/R_Y=Y'
\]
\textit{是 \(X'=X/R\) 的一个子对象, 且 \(Y'\to X'\) 属于 \((N)\).}

\textit{映射 \(Y\mapsto Y'\) 是一个双射,
其定义域为 \(X\) 的那些在 \(R\) 下稳定且使 \(Y\to X\)
属于 \((N)\) 的子对象 \(Y\) 所成的集合,
值域为 \(X'\) 的那些使 \(Y'\to X'\) 属于 \((N)\)
的子对象 \(Y'\) 所成的集合. 其逆映射为}
\[
  Y'\longmapsto Y'\times_{X'}X.
\]

事实上, 由于 \(R\) 是 \((M)\)-有效的, 态射 \(X\to X'\) 属于 \((M)\).
设 \(Y'\) 是 \(X'\) 的一个子对象,
使得典范态射 \(Y'\to X'\) 属于 \((N)\).
则 \(Y=Y'\times_{X'}X\) 是 \(X\) 的一个在 \(R\) 下稳定的子对象,
且态射 \(Y\to X\) 和 \(Y\to Y'\) 分别属于 \((N)\) 和 \((M)\),
因为这两个族在基变换下都稳定.
以 \(R_Y\) 表示 \(R\) 在 \(Y\) 中诱导的等价关系.
由 \SGARefLink{sga3:iv:proposition:4:4:11}{4.4.11},
层商 \(Y/R_Y\) 由 \(Y'\) 表出;
因此 \SGARefLink{sga3:iv:lemma:4:6:5}{4.6.5} 表明 \(R_Y\) 是 \((M)\)-有效的.

反之, 设 \(Y\) 是 \(X\) 的一个在 \(R\) 下稳定的子对象,
并且 \(Y\to X\) 属于 \((N)\).
由于 \(Y\) 在 \(R\) 下稳定,
它在
\[
  R=X\times_{X'}X
\]
中的两个逆像相同, 从而 \(Y\) 配备了关于 \(X\to X'\) 的一个下降数据.
由族 \((N)\) 的公理
\SGARefLink{sga3:iv:subsection:4:7:item:f:m:01}{\((f_M)\)} 即得所需结论.

\medskip
\noindent\SGARefTarget{sga3:iv:corollary:4:7:3}\textbf{推论 4.7.3.}
\textit{设 \(X\in\operatorname{Ob}(\mathscr C)\),
\(R\) 是 \(X\) 中的一个 \((M)\)-有效等价关系.
再假设 \(R\to X\times X\) 属于 \((N)\),
其中 \((N)\) 满足 \SGARefLink{sga3:iv:subsection:4:7:item:a:01}{(a)} 和
\SGARefLink{sga3:iv:subsection:4:7:item:f:mathcal:t:01}{\((f_{\mathcal T})\)}.
则对 \SGARefLink{sga3:iv:proposition:4:7:2}{4.7.2} 中的每个 \(Y\),
\(R_Y\to Y\times Y\) 也属于 \((N)\),
因而由 \SGARefLink{sga3:iv:proposition:4:7:1}{4.7.1} 得到}
\[
\begin{aligned}
 (Y/R_Y)(S)=\{Z\subseteq Y_S\mid{}&
   Z\text{ 是一个 }S\text{-子对象, 且在 }R_Y\times S\text{ 下稳定},\\
 & Z\to Y_S\text{ 属于 }(N)
   \text{ (所以 }Z\to X_S\text{ 也属于该族}),\\
 & Z\to S\text{ 是覆盖态射},\quad R_Z\simeq Z\times_S Z\,\}.
\end{aligned}
\]

% Translation coverage: entire assigned source component, without omissions.
% Statistics: source lines=124; target lines=142; inline math=120; display math=5; unresolved translation TODOs=0.
